\paragraph{Functional Logic Programming Languages}

The combination of functional and logic programming has a long history, surveyed
by
\textcite{bellia_1986,darlington_1986,hanus_1994,antoy_2010}. The most
prominent FLP language is Curry \parencite{hanus_2013,hanus_2016}, which extends
Haskell with logical variables, equational constraints, and nondeterministic
functions. Curry has multiple implementations, including KiCS2
\parencite{brassel_2011}, which compiles Curry to Haskell, and PAKCS
\parencite{hanus_2022}, which targets Prolog. \Gweek{} is intentionally much
simpler than Curry---it lacks type classes, modules, and I/O---but it makes the
underlying algebraic theory explicit in its design.

The Verse language \parencite{EpicGames_2023,augustsson_2023} is a recent FLP
language whose core calculus is an untyped $\lambda$-calculus with choice and
logical variables. Verse uses a stream semantics and supports encapsulated
search. \Gweek{} differs from Verse in being typed, in using CBPV as its
foundation, and in not supporting encapsulated search. As discussed in
\parencite{jones_2026}, encapsulated search requires moving beyond algebraic
effects to \emph{effect handlers}, and raises questions about the interaction of
encapsulation with the commutativity and idempotence of choice.

The miniKanren family of relational programming languages
\parencite{friedman_2005,friedman_2018,byrd_2009} embeds logic programming in a
host language (typically Scheme). miniKanren is conceptually close to \Gweek{} in
its use of logical variables and unification, but it is purely relational: it
lacks the functional programming features (higher-order functions, algebraic
data types, pattern matching) that \Gweek{} inherits from CBPV.

\paragraph{Semantics of FLP}

The semantic foundations of \Gweek{} are developed in our companion paper
\parencite{jones_2026}. That paper presents a domain-theoretic semantics for a
CBPV-based FLP calculus, based on the lower powerdomain, and proves it sound,
adequate, and fully abstract for a termination-insensitive notion of
observational equivalence. The equational theory derived from this semantics
validates the rewrite rules used in \Gweek{}'s peephole optimiser.

Previous semantic work on FLP includes the CRWL rewriting logic framework of
\textcite{gonzalez-moreno_1999}, the operational semantics of
\textcite{albert_2005}, and the denotational approaches of
\textcite{christiansen_2011} and \textcite{mehner_2014}. The monad-based
approach to FLP effects, which \Gweek{} follows, originates with
\textcite{staton_2013}. The observation that logical variables are an effect that
can be added to a language in isolation goes back to \textcite{lindstrom_1985}.

\paragraph{Narrowing}

Narrowing \parencite{reddy_1985} is the standard execution strategy for FLP.
\Gweek{} implements \emph{weak head narrowing}: a logical variable is refined
only when it is scrutinised by a case expression or required by unification.
This is related to the \emph{needed narrowing} strategy of
\textcite{antoy_1997}, which refines variables only when they are `needed' by
the outermost function. In the CBPV framework, the distinction between values
and computations provides a clean account of when narrowing must occur: it
happens precisely when a computation attempts to eliminate a value that turns out
to be a logic variable. The completeness of this strategy is proved in
\parencite{jones_2026}.

\paragraph{Implementation Techniques}

\Gweek{}'s use of arena allocation for term representation is standard in
high-performance language implementations; see e.g.\ the Zig compiler and the
Rust compiler's use of arenas. The persistent cons-list representation of
environments is a classical technique in functional language implementations
\parencite{landin_1964,henderson_1980}. The union-find structure for logic
variables is the standard choice in logic programming implementations
\parencite{warren_1983}.

Embeddings of logic programming in functional languages
\parencite{claessen_2001,seres_1999,hinze_1998} provide similar functionality
to \Gweek{} but at the cost of embedding overhead and the loss of the host
language's type system for the logical fragment. As a standalone language,
\Gweek{} retains control over term representation and search implementation.
